% Options for packages loaded elsewhere
\PassOptionsToPackage{unicode}{hyperref}
\PassOptionsToPackage{hyphens}{url}
\documentclass[
  11pt,
  a4paper,
]{article}
\usepackage{xcolor}
\usepackage[margin=18mm]{geometry}
\usepackage{amsmath,amssymb}
\setcounter{secnumdepth}{-\maxdimen} % remove section numbering
\usepackage{iftex}
\ifPDFTeX
  \usepackage[T1]{fontenc}
  \usepackage[utf8]{inputenc}
  \usepackage{textcomp} % provide euro and other symbols
\else % if luatex or xetex
  \usepackage{unicode-math} % this also loads fontspec
  \defaultfontfeatures{Scale=MatchLowercase}
  \defaultfontfeatures[\rmfamily]{Ligatures=TeX,Scale=1}
\fi
\ifPDFTeX\else
  % xetex/luatex font selection
    \setmainfont[]{Noto Serif}
    \setsansfont[]{Noto Sans}
    \setmonofont[]{Noto Sans Mono}
  \setmathfont[]{Noto Sans Math}
\fi
% Use upquote if available, for straight quotes in verbatim environments
\IfFileExists{upquote.sty}{\usepackage{upquote}}{}
\IfFileExists{microtype.sty}{% use microtype if available
  \usepackage[]{microtype}
  \UseMicrotypeSet[protrusion]{basicmath} % disable protrusion for tt fonts
}{}
\makeatletter
\@ifundefined{KOMAClassName}{% if non-KOMA class
  \IfFileExists{parskip.sty}{%
    \usepackage{parskip}
  }{% else
    \setlength{\parindent}{0pt}
    \setlength{\parskip}{6pt plus 2pt minus 1pt}}
}{% if KOMA class
  \KOMAoptions{parskip=half}}
\makeatother
\ifLuaTeX
\usepackage[bidi=basic]{babel}
\else
\usepackage[bidi=default]{babel}
\fi
\babelprovide[main,import]{russian}
\ifPDFTeX
\else
\babelfont{rm}[]{Noto Serif}
\fi
% get rid of language-specific shorthands (see #6817):
\let\LanguageShortHands\languageshorthands
\def\languageshorthands#1{}
\setlength{\emergencystretch}{3em} % prevent overfull lines
\providecommand{\tightlist}{%
  \setlength{\itemsep}{0pt}\setlength{\parskip}{0pt}}
\usepackage{bookmark}
\IfFileExists{xurl.sty}{\usepackage{xurl}}{} % add URL line breaks if available
\urlstyle{same}
\hypersetup{
  pdftitle={Билеты по дискретным случайным процессам},
  pdflang={ru-RU},
  hidelinks,
  pdfcreator={LaTeX via pandoc}}

\title{Билеты по дискретным случайным процессам}
\author{}
\date{}

\begin{document}
\maketitle

\section{Билет 07. Простые функции, интеграл Лебега и теоремы Леви и
Лебега о предельном
переходе}\label{ux431ux438ux43bux435ux442-07.-ux43fux440ux43eux441ux442ux44bux435-ux444ux443ux43dux43aux446ux438ux438-ux438ux43dux442ux435ux433ux440ux430ux43b-ux43bux435ux431ux435ux433ux430-ux438-ux442ux435ux43eux440ux435ux43cux44b-ux43bux435ux432ux438-ux438-ux43bux435ux431ux435ux433ux430-ux43e-ux43fux440ux435ux434ux435ux43bux44cux43dux43eux43c-ux43fux435ux440ux435ux445ux43eux434ux435}

Источники: Ширяев 1, гл. II, §6; лекции ДСП.

\subsection{1. Короткий
ответ}\label{ux43aux43eux440ux43eux442ux43aux438ux439-ux43eux442ux432ux435ux442}

Интеграл Лебега по неотрицательной мере строится в три шага: сначала для
индикаторов, затем для простых неотрицательных функций, затем для
произвольных неотрицательных измеримых функций через супремум по простым
минорантам.

Если

\[
s=\sum_{k=1}^n a_k\mathbf 1_{A_k},
\qquad
a_k\ge0,\quad A_k\in\mathcal F,
\]

и множества \(A_k\) попарно не пересекаются, то

\[
\int s\,d\mu=\sum_{k=1}^n a_k\mu(A_k).
\]

Для неотрицательной измеримой функции \(f\):

\[
\int f\,d\mu
=
\sup\left\{\int s\,d\mu:\ 0\le s\le f,\ s\text{ простая}\right\}.
\]

Теорема Беппо Леви: если \(0\le f_n\uparrow f\), то

\[
\int f_n\,d\mu\uparrow\int f\,d\mu.
\]

Мажорантная теорема Лебега: если \(f_n\to f\) почти всюду и
\(|f_n|\le g\), где \(g\in L^1\), то

\[
\int f_n\,d\mu\to\int f\,d\mu.
\]

\subsection{2. Полный
ответ}\label{ux43fux43eux43bux43dux44bux439-ux43eux442ux432ux435ux442}

\subsubsection{2.1. Введение и роль интеграла
Лебега}\label{ux432ux432ux435ux434ux435ux43dux438ux435-ux438-ux440ux43eux43bux44c-ux438ux43dux442ux435ux433ux440ux430ux43bux430-ux43bux435ux431ux435ux433ux430}

В предыдущих билетах были введены измеримые функции и распределения.
Теперь нужен аппарат, который позволяет интегрировать измеримые функции
по мере. В вероятности этот интеграл становится математическим
ожиданием, а в дальнейшем через него выражаются моменты,
характеристические функции, условные ожидания, корреляции и спектральные
характеристики.

\subsubsection{2.2. Измеримое пространство и интеграл от индикаторной
функции}\label{ux438ux437ux43cux435ux440ux438ux43cux43eux435-ux43fux440ux43eux441ux442ux440ux430ux43dux441ux442ux432ux43e-ux438-ux438ux43dux442ux435ux433ux440ux430ux43b-ux43eux442-ux438ux43dux434ux438ux43aux430ux442ux43eux440ux43dux43eux439-ux444ux443ux43dux43aux446ux438ux438}

Пусть задано пространство с неотрицательной мерой

\[
(\Omega,\mathcal F,\mu).
\]

Неотрицательность меры означает, что

\[
\mu(A)\in[0,\infty]
\qquad
A\in\mathcal F.
\]

Сначала интеграл определяют для самой простой измеримой функции -
индикатора множества:

\[
\int \mathbf 1_A\,d\mu=\mu(A).
\]

Это естественно: индикатор равен \(1\) на \(A\) и \(0\) вне \(A\),
значит площадь под ним должна равняться мере множества \(A\).

\subsubsection{2.3. Интеграл от простых
функций}\label{ux438ux43dux442ux435ux433ux440ux430ux43b-ux43eux442-ux43fux440ux43eux441ux442ux44bux445-ux444ux443ux43dux43aux446ux438ux439}

Дальше рассматривают простые функции. Простая функция - измеримая
функция, принимающая конечное число значений. Если \(s\ge0\) и

\[
s=\sum_{k=1}^n a_k\mathbf 1_{A_k},
\]

где \(a_k\ge0\), \(A_k\in\mathcal F\) и \(A_k\) попарно не пересекаются,
то интеграл задают формулой

\[
\int s\,d\mu=\sum_{k=1}^n a_k\mu(A_k).
\]

Эта формула не зависит от выбранного представления простой функции: если
разложение измельчить до общего конечного разбиения, сумма останется той
же. Поэтому интеграл простой функции определен корректно.

\subsubsection{2.4. Построение интеграла для неотрицательных измеримых
функций}\label{ux43fux43eux441ux442ux440ux43eux435ux43dux438ux435-ux438ux43dux442ux435ux433ux440ux430ux43bux430-ux434ux43bux44f-ux43dux435ux43eux442ux440ux438ux446ux430ux442ux435ux43bux44cux43dux44bux445-ux438ux437ux43cux435ux440ux438ux43cux44bux445-ux444ux443ux43dux43aux446ux438ux439}

Для произвольной неотрицательной измеримой функции \(f\) интеграл строят
приближением снизу простыми функциями:

\[
\int f\,d\mu
=
\sup\left\{\int s\,d\mu:\ 0\le s\le f,\ s\text{ простая измеримая}\right\}.
\]

Интуиция такая: сложную функцию можно снизу аппроксимировать
ступенчатыми измеримыми функциями. Интеграл \(f\) - это наибольший
предел площадей таких нижних ступенчатых приближений. Значение интеграла
может быть бесконечным.

\subsubsection{2.5. Функции произвольного знака и условие
интегрируемости}\label{ux444ux443ux43dux43aux446ux438ux438-ux43fux440ux43eux438ux437ux432ux43eux43bux44cux43dux43eux433ux43e-ux437ux43dux430ux43aux430-ux438-ux443ux441ux43bux43eux432ux438ux435-ux438ux43dux442ux435ux433ux440ux438ux440ux443ux435ux43cux43eux441ux442ux438}

Если функция \(f\) может принимать отрицательные значения, ее
раскладывают на положительную и отрицательную части:

\[
f^+=\max(f,0),\qquad
f^-=\max(-f,0),
\]

так что

\[
f=f^+-f^-,
\qquad
|f|=f^++f^-.
\]

Если хотя бы один из интегралов \(\int f^+\,d\mu\), \(\int f^-\,d\mu\)
конечен, можно определить расширенный интеграл

\[
\int f\,d\mu=\int f^+\,d\mu-\int f^-\,d\mu,
\]

не допуская неопределенности \(\infty-\infty\). Если

\[
\int |f|\,d\mu<\infty,
\]

то \(f\) называется интегрируемой, или суммируемой, и пишут
\(f\in L^1(\mu)\).

\subsubsection{2.6. Теорема Беппо Леви о монотонной
сходимости}\label{ux442ux435ux43eux440ux435ux43cux430-ux431ux435ux43fux43fux43e-ux43bux435ux432ux438-ux43e-ux43cux43eux43dux43eux442ux43eux43dux43dux43eux439-ux441ux445ux43eux434ux438ux43cux43eux441ux442ux438}

Главная сила интеграла Лебега - хорошие теоремы о предельном переходе. В
программе эти теоремы идут без доказательства, но условия нужно знать
точно.

Первая - теорема Беппо Леви, или теорема о монотонной сходимости. Если

\[
0\le f_1\le f_2\le\cdots
\]

и

\[
f_n(\omega)\uparrow f(\omega)
\]

для всех \(\omega\) или почти всюду, то

\[
\int f_n\,d\mu\uparrow\int f\,d\mu.
\]

Здесь не нужен интегрируемый мажорант; зато нужна неотрицательность и
монотонный рост.

\subsubsection{2.7. Мажорантная теорема Лебега о доминированной
сходимости}\label{ux43cux430ux436ux43eux440ux430ux43dux442ux43dux430ux44f-ux442ux435ux43eux440ux435ux43cux430-ux43bux435ux431ux435ux433ux430-ux43e-ux434ux43eux43cux438ux43dux438ux440ux43eux432ux430ux43dux43dux43eux439-ux441ux445ux43eux434ux438ux43cux43eux441ux442ux438}

Вторая - мажорантная теорема Лебега, или теорема Лебега о доминированной
сходимости. Если

\[
f_n\to f
\]

почти всюду и существует интегрируемая функция \(g\in L^1(\mu)\) такая,
что

\[
|f_n|\le g
\qquad \text{почти всюду для всех }n,
\]

то \(f\) интегрируема и

\[
\int f_n\,d\mu\to\int f\,d\mu.
\]

Эта теорема позволяет менять предел и интеграл при наличии общего
интегрируемого контроля. Без такого контроля утверждение может быть
неверным.

\subsection{3. Основные
определения}\label{ux43eux441ux43dux43eux432ux43dux44bux435-ux43eux43fux440ux435ux434ux435ux43bux435ux43dux438ux44f}

\subsubsection{Пространство с неотрицательной
мерой}\label{ux43fux440ux43eux441ux442ux440ux430ux43dux441ux442ux432ux43e-ux441-ux43dux435ux43eux442ux440ux438ux446ux430ux442ux435ux43bux44cux43dux43eux439-ux43cux435ux440ux43eux439}

Тройка \((\Omega,\mathcal F,\mu)\), где \(\mathcal F\) -
\(\sigma\)-алгебра, а \(\mu:\mathcal F\to[0,\infty]\) -
счетно-аддитивная неотрицательная мера.

\subsubsection{Индикатор
множества}\label{ux438ux43dux434ux438ux43aux430ux442ux43eux440-ux43cux43dux43eux436ux435ux441ux442ux432ux430}

Для \(A\in\mathcal F\):

\[
\mathbf 1_A(\omega)=
\begin{cases}
1, & \omega\in A,\\
0, & \omega\notin A.
\end{cases}
\]

\subsubsection{Простая
функция}\label{ux43fux440ux43eux441ux442ux430ux44f-ux444ux443ux43dux43aux446ux438ux44f}

Измеримая функция \(s:\Omega\to\mathbb R\) называется простой, если она
принимает конечное число значений. Ее можно записать в виде

\[
s=\sum_{k=1}^n a_k\mathbf 1_{A_k},
\]

где \(A_k\in\mathcal F\). В каноническом представлении значения \(a_k\)
различны, а множества

\[
A_k=\{\omega:s(\omega)=a_k\}
\]

попарно не пересекаются и образуют разбиение области, где \(s\)
принимает эти значения.

\subsubsection{Неотрицательная простая
функция}\label{ux43dux435ux43eux442ux440ux438ux446ux430ux442ux435ux43bux44cux43dux430ux44f-ux43fux440ux43eux441ux442ux430ux44f-ux444ux443ux43dux43aux446ux438ux44f}

Простая функция \(s\), для которой \(s(\omega)\ge0\) для всех
\(\omega\). Если неравенство выполнено только почти всюду, функцию можно
заменить равной ей почти всюду неотрицательной модификацией.

\subsubsection{Интеграл
индикатора}\label{ux438ux43dux442ux435ux433ux440ux430ux43b-ux438ux43dux434ux438ux43aux430ux442ux43eux440ux430}

\[
\int \mathbf 1_A\,d\mu=\mu(A).
\]

\subsubsection{Интеграл неотрицательной простой
функции}\label{ux438ux43dux442ux435ux433ux440ux430ux43b-ux43dux435ux43eux442ux440ux438ux446ux430ux442ux435ux43bux44cux43dux43eux439-ux43fux440ux43eux441ux442ux43eux439-ux444ux443ux43dux43aux446ux438ux438}

Если

\[
s=\sum_{k=1}^n a_k\mathbf 1_{A_k},
\qquad
a_k\ge0,
\]

и \(A_k\) попарно не пересекаются, то

\[
\int s\,d\mu=\sum_{k=1}^n a_k\mu(A_k).
\]

\subsubsection{Интеграл неотрицательной измеримой
функции}\label{ux438ux43dux442ux435ux433ux440ux430ux43b-ux43dux435ux43eux442ux440ux438ux446ux430ux442ux435ux43bux44cux43dux43eux439-ux438ux437ux43cux435ux440ux438ux43cux43eux439-ux444ux443ux43dux43aux446ux438ux438}

Для \(f:\Omega\to[0,\infty]\):

\[
\int f\,d\mu
=
\sup\left\{\int s\,d\mu:\ 0\le s\le f,\ s\text{ простая измеримая}\right\}.
\]

\subsubsection{Положительная и отрицательная части
функции}\label{ux43fux43eux43bux43eux436ux438ux442ux435ux43bux44cux43dux430ux44f-ux438-ux43eux442ux440ux438ux446ux430ux442ux435ux43bux44cux43dux430ux44f-ux447ux430ux441ux442ux438-ux444ux443ux43dux43aux446ux438ux438}

\[
f^+=\max(f,0),\qquad
f^-=\max(-f,0).
\]

Тогда

\[
f=f^+-f^-,
\qquad
|f|=f^++f^-.
\]

\subsubsection{Интегрируемая
функция}\label{ux438ux43dux442ux435ux433ux440ux438ux440ux443ux435ux43cux430ux44f-ux444ux443ux43dux43aux446ux438ux44f}

Измеримая функция \(f\) называется интегрируемой, если

\[
\int |f|\,d\mu<\infty.
\]

В этом случае

\[
\int f\,d\mu=\int f^+\,d\mu-\int f^-\,d\mu
\]

и оба интеграла справа конечны.

\subsubsection{Почти
всюду}\label{ux43fux43eux447ux442ux438-ux432ux441ux44eux434ux443}

Свойство выполнено почти всюду, если оно нарушается только на множестве
меры ноль.

\subsubsection{Мажоранта}\label{ux43cux430ux436ux43eux440ux430ux43dux442ux430}

Функция \(g\) называется интегрируемой мажорантой для последовательности
\(f_n\), если

\[
|f_n|\le g\quad \text{почти всюду для всех }n,
\qquad
g\in L^1(\mu).
\]

\subsection{4. Основные теоремы и
свойства}\label{ux43eux441ux43dux43eux432ux43dux44bux435-ux442ux435ux43eux440ux435ux43cux44b-ux438-ux441ux432ux43eux439ux441ux442ux432ux430}

\subsubsection{Свойство 1. Корректность интеграла простой
функции}\label{ux441ux432ux43eux439ux441ux442ux432ux43e-1.-ux43aux43eux440ux440ux435ux43aux442ux43dux43eux441ux442ux44c-ux438ux43dux442ux435ux433ux440ux430ux43bux430-ux43fux440ux43eux441ux442ux43eux439-ux444ux443ux43dux43aux446ux438ux438}

Если неотрицательная простая функция имеет два представления

\[
s=\sum_{i=1}^m a_i\mathbf 1_{A_i}
=
\sum_{j=1}^n b_j\mathbf 1_{B_j},
\]

где множества в каждом представлении попарно не пересекаются, то

\[
\sum_{i=1}^m a_i\mu(A_i)
=
\sum_{j=1}^n b_j\mu(B_j).
\]

\subsubsection{Доказательство-идея}\label{ux434ux43eux43aux430ux437ux430ux442ux435ux43bux44cux441ux442ux432ux43e-ux438ux434ux435ux44f}

Перейдем к общему измельчению

\[
C_{ij}=A_i\cap B_j.
\]

На каждом непустом \(C_{ij}\) значения двух представлений совпадают, то
есть \(a_i=b_j\). Поэтому обе суммы после разбиения на \(C_{ij}\)
становятся одной и той же суммой.

\subsubsection{Свойство 2. Интеграл индикатора равен
мере}\label{ux441ux432ux43eux439ux441ux442ux432ux43e-2.-ux438ux43dux442ux435ux433ux440ux430ux43b-ux438ux43dux434ux438ux43aux430ux442ux43eux440ux430-ux440ux430ux432ux435ux43d-ux43cux435ux440ux435}

Для любого \(A\in\mathcal F\):

\[
\int \mathbf 1_A\,d\mu=\mu(A).
\]

Это базовая нормировка интеграла Лебега.

\subsubsection{Свойство 3. Монотонность
интеграла}\label{ux441ux432ux43eux439ux441ux442ux432ux43e-3.-ux43cux43eux43dux43eux442ux43eux43dux43dux43eux441ux442ux44c-ux438ux43dux442ux435ux433ux440ux430ux43bux430}

Если \(0\le f\le g\), то

\[
\int f\,d\mu\le\int g\,d\mu.
\]

\subsubsection{Доказательство}\label{ux434ux43eux43aux430ux437ux430ux442ux435ux43bux44cux441ux442ux432ux43e}

Любая простая функция \(s\), удовлетворяющая \(0\le s\le f\),
автоматически удовлетворяет \(0\le s\le g\). Поэтому супремум по простым
минорантам для \(f\) не больше супремума для \(g\).

\subsubsection{Свойство 4. Линейность для неотрицательных
функций}\label{ux441ux432ux43eux439ux441ux442ux432ux43e-4.-ux43bux438ux43dux435ux439ux43dux43eux441ux442ux44c-ux434ux43bux44f-ux43dux435ux43eux442ux440ux438ux446ux430ux442ux435ux43bux44cux43dux44bux445-ux444ux443ux43dux43aux446ux438ux439}

Если \(f,g\ge0\) измеримы и \(a,b\ge0\), то

\[
\int(af+bg)\,d\mu
=
a\int f\,d\mu+b\int g\,d\mu.
\]

Здесь допускаются бесконечные значения, если не возникает противоречивых
операций.

\subsubsection{Доказательство-идея}\label{ux434ux43eux43aux430ux437ux430ux442ux435ux43bux44cux441ux442ux432ux43e-ux438ux434ux435ux44f-1}

Для простых функций это следует из формулы суммы по разбиению. Для общих
неотрицательных функций берут возрастающие последовательности простых
функций, приближающие \(f\) и \(g\) снизу, и применяют теорему Беппо
Леви.

\subsubsection{\texorpdfstring{Свойство 5. Линейность на
\(L^1\)}{Свойство 5. Линейность на L\^{}1}}\label{ux441ux432ux43eux439ux441ux442ux432ux43e-5.-ux43bux438ux43dux435ux439ux43dux43eux441ux442ux44c-ux43dux430-l1}

Если \(f,g\in L^1(\mu)\) и \(a,b\in\mathbb R\), то

\[
\int(af+bg)\,d\mu
=
a\int f\,d\mu+b\int g\,d\mu.
\]

\subsubsection{Обоснование}\label{ux43eux431ux43eux441ux43dux43eux432ux430ux43dux438ux435}

Интегрируемость гарантирует конечность интегралов положительных и
отрицательных частей, поэтому не возникает неопределенности
\(\infty-\infty\).

\subsubsection{Свойство 6. Если функция равна нулю почти всюду, то ее
интеграл равен
нулю}\label{ux441ux432ux43eux439ux441ux442ux432ux43e-6.-ux435ux441ux43bux438-ux444ux443ux43dux43aux446ux438ux44f-ux440ux430ux432ux43dux430-ux43dux443ux43bux44e-ux43fux43eux447ux442ux438-ux432ux441ux44eux434ux443-ux442ux43e-ux435ux435-ux438ux43dux442ux435ux433ux440ux430ux43b-ux440ux430ux432ux435ux43d-ux43dux443ux43bux44e}

Если \(f\ge0\) и \(f=0\) почти всюду, то

\[
\int f\,d\mu=0.
\]

Если \(f=g\) почти всюду и обе функции интегрируемы, то

\[
\int f\,d\mu=\int g\,d\mu.
\]

\subsubsection{Свойство 7. Аппроксимация неотрицательной измеримой
функции
простыми}\label{ux441ux432ux43eux439ux441ux442ux432ux43e-7.-ux430ux43fux43fux440ux43eux43aux441ux438ux43cux430ux446ux438ux44f-ux43dux435ux43eux442ux440ux438ux446ux430ux442ux435ux43bux44cux43dux43eux439-ux438ux437ux43cux435ux440ux438ux43cux43eux439-ux444ux443ux43dux43aux446ux438ux438-ux43fux440ux43eux441ux442ux44bux43cux438}

Для любой неотрицательной измеримой функции \(f\) существует
последовательность неотрицательных простых функций \(s_n\) такая, что

\[
0\le s_1\le s_2\le\cdots\le f,
\qquad
s_n(\omega)\uparrow f(\omega).
\]

\subsubsection{Доказательство-идея}\label{ux434ux43eux43aux430ux437ux430ux442ux435ux43bux44cux441ux442ux432ux43e-ux438ux434ux435ux44f-2}

Можно взять ступенчатое округление вниз. Например, на уровне \(n\)
разбить значения \(f\) на интервалы длины \(2^{-n}\) до высоты \(n\) и
положить

\[
s_n(\omega)=\frac{k}{2^n}
\]

на множестве

\[
\left\{\frac{k}{2^n}\le f(\omega)<\frac{k+1}{2^n}\right\},
\]

а для \(f(\omega)\ge n\) ограничить значение сверху числом \(n\). Эти
множества измеримы, потому что \(f\) измерима.

\subsubsection{Теорема 8. Беппо Леви, или монотонная
сходимость}\label{ux442ux435ux43eux440ux435ux43cux430-8.-ux431ux435ux43fux43fux43e-ux43bux435ux432ux438-ux438ux43bux438-ux43cux43eux43dux43eux442ux43eux43dux43dux430ux44f-ux441ux445ux43eux434ux438ux43cux43eux441ux442ux44c}

Пусть \(f_n:\Omega\to[0,\infty]\) измеримы,

\[
0\le f_1\le f_2\le\cdots,
\]

и

\[
f_n(\omega)\uparrow f(\omega)
\]

почти всюду. Тогда

\[
\int f_n\,d\mu\uparrow\int f\,d\mu.
\]

В программе эта теорема идет без доказательства.

\subsubsection{Идея
доказательства}\label{ux438ux434ux435ux44f-ux434ux43eux43aux430ux437ux430ux442ux435ux43bux44cux441ux442ux432ux430}

Неравенство

\[
\lim_n\int f_n\,d\mu\le\int f\,d\mu
\]

следует из монотонности. Обратное неравенство получают так: берут
простую функцию \(s\le f\), немного уменьшают ее до \(\alpha s\), где
\(0<\alpha<1\), и используют множества, на которых \(f_n\) уже превысила
\(\alpha s\). Потом переходят к пределу по \(n\) и по
\(\alpha\uparrow1\).

\subsubsection{Теорема 9. Мажорантная теорема Лебега, или доминированная
сходимость}\label{ux442ux435ux43eux440ux435ux43cux430-9.-ux43cux430ux436ux43eux440ux430ux43dux442ux43dux430ux44f-ux442ux435ux43eux440ux435ux43cux430-ux43bux435ux431ux435ux433ux430-ux438ux43bux438-ux434ux43eux43cux438ux43dux438ux440ux43eux432ux430ux43dux43dux430ux44f-ux441ux445ux43eux434ux438ux43cux43eux441ux442ux44c}

Пусть \(f_n\) измеримы,

\[
f_n\to f
\]

почти всюду, и существует \(g\in L^1(\mu)\) такая, что

\[
|f_n|\le g
\qquad \text{почти всюду для всех }n.
\]

Тогда \(f\in L^1(\mu)\) и

\[
\int f_n\,d\mu\to\int f\,d\mu.
\]

Более сильная форма вывода:

\[
\int |f_n-f|\,d\mu\to0.
\]

Из нее уже следует сходимость интегралов.

В программе эта теорема также идет без полного доказательства.

\subsubsection{Идея
доказательства}\label{ux438ux434ux435ux44f-ux434ux43eux43aux430ux437ux430ux442ux435ux43bux44cux441ux442ux432ux430-1}

Из \(|f_n|\le g\) и \(f_n\to f\) следует \(|f|\le g\), значит \(f\)
интегрируема. Далее к неотрицательным функциям

\[
2g-|f_n-f|
\]

или к верхним и нижним оценкам применяют лемму Фату. Смысл:
интегрируемая мажоранта запрещает массе убегать в бесконечность или
концентрироваться на все меньших множествах с большой высотой.

\subsubsection{Свойство 10. Нужны именно условия предельных
теорем}\label{ux441ux432ux43eux439ux441ux442ux432ux43e-10.-ux43dux443ux436ux43dux44b-ux438ux43cux435ux43dux43dux43e-ux443ux441ux43bux43eux432ux438ux44f-ux43fux440ux435ux434ux435ux43bux44cux43dux44bux445-ux442ux435ux43eux440ux435ux43c}

Если \(f_n\to f\) почти всюду, но нет ни монотонности, ни интегрируемой
мажоранты, то из этого не следует

\[
\int f_n\,d\mu\to\int f\,d\mu.
\]

Например, на \([0,1]\) с мерой Лебега функции

\[
f_n(x)=n\mathbf 1_{(0,1/n)}(x)
\]

сходятся к \(0\) почти всюду, но

\[
\int_0^1 f_n(x)\,dx=1
\]

для всех \(n\).

\subsection{5. Примеры}\label{ux43fux440ux438ux43cux435ux440ux44b}

\subsubsection{Пример 1.
Индикатор}\label{ux43fux440ux438ux43cux435ux440-1.-ux438ux43dux434ux438ux43aux430ux442ux43eux440}

Если \(A\in\mathcal F\), то

\[
\int \mathbf 1_A\,d\mu=\mu(A).
\]

В вероятностном пространстве \((\Omega,\mathcal F,P)\):

\[
\mathbb E\mathbf 1_A=P(A).
\]

\subsubsection{Пример 2. Простая
функция}\label{ux43fux440ux438ux43cux435ux440-2.-ux43fux440ux43eux441ux442ux430ux44f-ux444ux443ux43dux43aux446ux438ux44f}

Пусть \(A,B\in\mathcal F\), \(A\cap B=\varnothing\), и

\[
s=2\mathbf 1_A+5\mathbf 1_B.
\]

Тогда

\[
\int s\,d\mu=2\mu(A)+5\mu(B).
\]

Если \(\mu(A)=3\), \(\mu(B)=1\), то

\[
\int s\,d\mu=11.
\]

\subsubsection{Пример 3. Дискретное математическое
ожидание}\label{ux43fux440ux438ux43cux435ux440-3.-ux434ux438ux441ux43aux440ux435ux442ux43dux43eux435-ux43cux430ux442ux435ux43cux430ux442ux438ux447ux435ux441ux43aux43eux435-ux43eux436ux438ux434ux430ux43dux438ux435}

Пусть \(X\) принимает значения \(x_1,\ldots,x_n\) с вероятностями
\(p_1,\ldots,p_n\). Тогда

\[
X=\sum_{k=1}^n x_k\mathbf 1_{\{X=x_k\}},
\]

и

\[
\mathbb E X=\int X\,dP=\sum_{k=1}^n x_kp_k
\]

при условии, что \(X\) неотрицательна или интегрируема. Это обычная
формула ожидания для дискретной величины.

\subsubsection{Пример 4. Интеграл по
плотности}\label{ux43fux440ux438ux43cux435ux440-4.-ux438ux43dux442ux435ux433ux440ux430ux43b-ux43fux43e-ux43fux43bux43eux442ux43dux43eux441ux442ux438}

Если \(X\) имеет плотность \(p\) на \(\mathbb R\), то вероятность
события \(X\in B\) выражается как интеграл индикатора по распределению:

\[
P\{X\in B\}=\int \mathbf 1_B(x)p(x)\,dx.
\]

Это связывает билет с распределениями из Билета 06.

\subsubsection{Пример 5. Теорема Беппо
Леви}\label{ux43fux440ux438ux43cux435ux440-5.-ux442ux435ux43eux440ux435ux43cux430-ux431ux435ux43fux43fux43e-ux43bux435ux432ux438}

На \([0,1]\) пусть

\[
f_n(x)=\mathbf 1_{[0,1-1/n]}(x).
\]

Тогда

\[
0\le f_n\uparrow \mathbf 1_{[0,1)}
\]

и

\[
\int_0^1 f_n(x)\,dx=1-\frac1n
\uparrow
1
=
\int_0^1 \mathbf 1_{[0,1)}(x)\,dx.
\]

\subsubsection{Пример 6. Мажорантная теорема
Лебега}\label{ux43fux440ux438ux43cux435ux440-6.-ux43cux430ux436ux43eux440ux430ux43dux442ux43dux430ux44f-ux442ux435ux43eux440ux435ux43cux430-ux43bux435ux431ux435ux433ux430}

На \([0,1]\) пусть

\[
f_n(x)=x^n.
\]

Тогда \(f_n(x)\to0\) для \(x\in[0,1)\) и \(f_n(1)=1\). Значит
\(f_n\to0\) почти всюду. Кроме того,

\[
0\le f_n(x)\le1,
\]

а \(g(x)=1\) интегрируема на \([0,1]\). Поэтому

\[
\int_0^1 x^n\,dx\to0.
\]

И действительно,

\[
\int_0^1 x^n\,dx=\frac1{n+1}.
\]

\subsubsection{Пример 7. Почему мажоранта
нужна}\label{ux43fux440ux438ux43cux435ux440-7.-ux43fux43eux447ux435ux43cux443-ux43cux430ux436ux43eux440ux430ux43dux442ux430-ux43dux443ux436ux43dux430}

На \([0,1]\) функции

\[
f_n(x)=n\mathbf 1_{(0,1/n)}(x)
\]

сходятся к \(0\) почти всюду, но

\[
\int_0^1 f_n\,dx=1.
\]

Значит без интегрируемой мажоранты нельзя просто переносить предел под
интеграл.

\subsection{6. Доказательства и
обоснования}\label{ux434ux43eux43aux430ux437ux430ux442ux435ux43bux44cux441ux442ux432ux430-ux438-ux43eux431ux43eux441ux43dux43eux432ux430ux43dux438ux44f}

\textbf{Почему простые функции являются первым шагом.}

Простая функция принимает конечное число значений, поэтому ее интеграл
сводится к конечной сумме мер уровней. Это ровно тот случай, где мера
множеств напрямую превращается в интеграл функции.

\textbf{Почему нужно приближение снизу.}

Для неотрицательной функции \(f\) нельзя сразу сложить значения по всем
уровням, потому что значений может быть бесконечно много. Вместо этого
берут все простые функции \(s\), которые лежат ниже \(f\), и выбирают
супремум их интегралов:

\[
\int f\,d\mu=\sup_{0\le s\le f}\int s\,d\mu.
\]

Так определение сохраняет монотонность и согласуется с интегралом
простых функций.

\textbf{Почему интеграл может быть бесконечным.}

Для неотрицательных функций значение \(\int f\,d\mu\) допускается в
\([0,\infty]\). Например, на \(\mathbb R\):

\[
\int_{\mathbb R}1\,dx=\infty.
\]

Это не ошибка, а часть конструкции. Конечность требуется только для
интегрируемости в смысле \(L^1\).

\textbf{Почему для знакопеременных функций нужна абсолютная
интегрируемость.}

Если обе части \(f^+\) и \(f^-\) имеют бесконечный интеграл, выражение

\[
\int f^+\,d\mu-\int f^-\,d\mu
\]

имеет вид \(\infty-\infty\) и не определено. Условие

\[
\int |f|\,d\mu<\infty
\]

гарантирует, что обе части конечны.

\textbf{Почему теорема Леви не требует мажоранты.}

В теореме Леви функции растут монотонно снизу. Интегралы не могут
колебаться или терять массу: они тоже образуют возрастающую
последовательность. Поэтому предел интегралов совпадает с интегралом
предельной функции, возможно бесконечным.

\textbf{Почему теорема Лебега требует мажоранту.}

При обычной почти всюду сходимости графики могут становиться очень
высокими на очень маленьких множествах. Точечный предел этого не видит,
а интеграл видит площадь. Интегрируемая мажоранта \(g\) запрещает такие
пики и дает контроль:

\[
|f_n|\le g,\qquad \int g\,d\mu<\infty.
\]

\textbf{Что значит ``без доказательства'' в программе.}

Для теорем Беппо Леви и Лебега на экзамене достаточно дать точные
условия и вывод, а также объяснить идею. Полные доказательства можно не
разворачивать, если преподаватель специально не попросит.

\subsection{7. Что написать на
доске}\label{ux447ux442ux43e-ux43dux430ux43fux438ux441ux430ux442ux44c-ux43dux430-ux434ux43eux441ux43aux435}

\begin{enumerate}
\def\labelenumi{\arabic{enumi}.}
\tightlist
\item
  Простая функция:
\end{enumerate}

\[
s=\sum_{k=1}^n a_k\mathbf 1_{A_k}.
\]

\begin{enumerate}
\def\labelenumi{\arabic{enumi}.}
\setcounter{enumi}{1}
\tightlist
\item
  Интеграл простой неотрицательной функции:
\end{enumerate}

\[
\int s\,d\mu=\sum_{k=1}^n a_k\mu(A_k).
\]

\begin{enumerate}
\def\labelenumi{\arabic{enumi}.}
\setcounter{enumi}{2}
\tightlist
\item
  Интеграл неотрицательной функции:
\end{enumerate}

\[
\int f\,d\mu
=
\sup_{0\le s\le f,\ s\text{ простая}}
\int s\,d\mu.
\]

\begin{enumerate}
\def\labelenumi{\arabic{enumi}.}
\setcounter{enumi}{3}
\tightlist
\item
  Положительная и отрицательная части:
\end{enumerate}

\[
f=f^+-f^-,
\qquad
|f|=f^++f^-.
\]

\begin{enumerate}
\def\labelenumi{\arabic{enumi}.}
\setcounter{enumi}{4}
\tightlist
\item
  Интегрируемость:
\end{enumerate}

\[
f\in L^1(\mu)
\quad\Longleftrightarrow\quad
\int |f|\,d\mu<\infty.
\]

\begin{enumerate}
\def\labelenumi{\arabic{enumi}.}
\setcounter{enumi}{5}
\tightlist
\item
  Беппо Леви:
\end{enumerate}

\[
0\le f_n\uparrow f
\quad\Longrightarrow\quad
\int f_n\,d\mu\uparrow\int f\,d\mu.
\]

\begin{enumerate}
\def\labelenumi{\arabic{enumi}.}
\setcounter{enumi}{6}
\tightlist
\item
  Лебег:
\end{enumerate}

\[
f_n\to f,\quad |f_n|\le g,\quad g\in L^1
\quad\Longrightarrow\quad
\int f_n\,d\mu\to\int f\,d\mu.
\]

\subsection{8. Типичные дополнительные
вопросы}\label{ux442ux438ux43fux438ux447ux43dux44bux435-ux434ux43eux43fux43eux43bux43dux438ux442ux435ux43bux44cux43dux44bux435-ux432ux43eux43fux440ux43eux441ux44b}

\textbf{Вопрос. Что такое простая функция?}

Это измеримая функция, принимающая конечное число значений. Ее можно
записать как

\[
s=\sum_{k=1}^n a_k\mathbf 1_{A_k},
\qquad A_k\in\mathcal F.
\]

\textbf{Вопрос. Почему интеграл индикатора равен мере множества?}

Индикатор \(\mathbf 1_A\) равен \(1\) на \(A\) и \(0\) вне \(A\),
поэтому площадь под ним равна размеру множества \(A\):

\[
\int\mathbf 1_A\,d\mu=\mu(A).
\]

\textbf{Вопрос. Как определяется интеграл неотрицательной измеримой
функции?}

Через супремум интегралов простых минорант:

\[
\int f\,d\mu=\sup_{0\le s\le f}\int s\,d\mu.
\]

\textbf{Вопрос. Что говорит теорема Беппо Леви?}

Если \(0\le f_n\uparrow f\), то

\[
\int f_n\,d\mu\uparrow\int f\,d\mu.
\]

\textbf{Вопрос. Чем теорема Леви отличается от теоремы Лебега?}

В теореме Леви нужна монотонность и неотрицательность, но не нужна
мажоранта. В теореме Лебега монотонность не нужна, зато нужна
интегрируемая мажоранта \(g\).

\textbf{Вопрос. Что такое интегрируемая мажоранта?}

Это функция \(g\in L^1(\mu)\) такая, что

\[
|f_n|\le g
\]

почти всюду для всех \(n\).

\textbf{Вопрос. Почему нельзя применять теорему Лебега без мажоранты?}

Почти всюду сходимость не контролирует интегралы. Пример:

\[
f_n=n\mathbf 1_{(0,1/n)}\to0
\]

почти всюду, но

\[
\int_0^1 f_n\,dx=1.
\]

\textbf{Вопрос. Что значит \(f\in L^1\)?}

Это значит, что \(f\) измерима и

\[
\int |f|\,d\mu<\infty.
\]

\textbf{Вопрос. Почему для знакопеременной функции нельзя просто
интегрировать положительную и отрицательную части без условий?}

Потому что может возникнуть неопределенность \(\infty-\infty\).
Абсолютная интегрируемость исключает эту проблему.

\subsection{9. Типичные
ошибки}\label{ux442ux438ux43fux438ux447ux43dux44bux435-ux43eux448ux438ux431ux43aux438}

\begin{enumerate}
\def\labelenumi{\arabic{enumi}.}
\item
  Применять мажорантную теорему Лебега, не предъявив \(g\in L^1\).
\item
  Путать теорему Беппо Леви и теорему Лебега: в Леви нужна монотонность,
  в Лебеге - мажоранта.
\item
  Забывать, что для неотрицательной функции интеграл может быть равен
  \(\infty\).
\item
  Для знакопеременной функции писать \(\int f=\int f^+-\int f^-\), не
  проверив отсутствие неопределенности \(\infty-\infty\).
\item
  Считать, что почти всюду сходимость сама по себе гарантирует
  сходимость интегралов.
\item
  Записывать интеграл простой функции без требования измеримости
  множеств \(A_k\).
\item
  Забывать, что представление простой функции можно измельчать, и
  интеграл не зависит от конкретной записи.
\end{enumerate}

\subsection{10. Связи с другими
билетами}\label{ux441ux432ux44fux437ux438-ux441-ux434ux440ux443ux433ux438ux43cux438-ux431ux438ux43bux435ux442ux430ux43cux438}

\textbf{Билет 03.} Интеграл Лебега строится по мере; мера Лебега
получается как продолжение меры с полуалгебры.

\textbf{Билет 04.} Для интеграла по \(\mathbb R^d\) важны борелевские и
лебегово измеримые множества.

\textbf{Билет 05.} Интегрировать можно измеримые функции; измеримость -
входное условие конструкции.

\textbf{Билет 06.} Вероятности через плотности записываются как
интегралы по мере Лебега.

\textbf{Билет 08.} Характеристические функции и интегрируемость
ограниченных измеримых функций опираются на интеграл Лебега.

\textbf{Билет 09.} Теоремы Фубини и независимость используют интегралы
по произведению мер.

\textbf{Билет 10.} Теорема Радона-Никодима и условное ожидание
формулируются через интегралы.

\textbf{Билет 11.} Замена переменной, моменты и дисперсии выражаются
через интеграл по вероятностному пространству и по распределению.

\subsection{11. Минимум на
удовлетворительно}\label{ux43cux438ux43dux438ux43cux443ux43c-ux43dux430-ux443ux434ux43eux432ux43bux435ux442ux432ux43eux440ux438ux442ux435ux43bux44cux43dux43e}

Нужно уметь сказать:

\begin{enumerate}
\def\labelenumi{\arabic{enumi}.}
\tightlist
\item
  Простая функция:
\end{enumerate}

\[
s=\sum_{k=1}^n a_k\mathbf 1_{A_k}.
\]

\begin{enumerate}
\def\labelenumi{\arabic{enumi}.}
\setcounter{enumi}{1}
\tightlist
\item
  Интеграл простой функции:
\end{enumerate}

\[
\int s\,d\mu=\sum_{k=1}^n a_k\mu(A_k).
\]

\begin{enumerate}
\def\labelenumi{\arabic{enumi}.}
\setcounter{enumi}{2}
\tightlist
\item
  Интеграл неотрицательной функции:
\end{enumerate}

\[
\int f\,d\mu=\sup_{0\le s\le f}\int s\,d\mu.
\]

\begin{enumerate}
\def\labelenumi{\arabic{enumi}.}
\setcounter{enumi}{3}
\tightlist
\item
  Теорема Беппо Леви:
\end{enumerate}

\[
0\le f_n\uparrow f
\Rightarrow
\int f_n\,d\mu\uparrow\int f\,d\mu.
\]

\begin{enumerate}
\def\labelenumi{\arabic{enumi}.}
\setcounter{enumi}{4}
\tightlist
\item
  Теорема Лебега:
\end{enumerate}

\[
f_n\to f,\quad |f_n|\le g\in L^1
\Rightarrow
\int f_n\,d\mu\to\int f\,d\mu.
\]

\subsection{12. Минимум на
хорошо}\label{ux43cux438ux43dux438ux43cux443ux43c-ux43dux430-ux445ux43eux440ux43eux448ux43e}

Кроме минимума на удовлетворительно, нужно:

\begin{enumerate}
\def\labelenumi{\arabic{enumi}.}
\tightlist
\item
  Объяснить, почему интеграл простой функции не зависит от
  представления.
\item
  Назвать монотонность и линейность интеграла.
\item
  Рассказать про положительную и отрицательную части функции.
\item
  Объяснить, что такое \(L^1\) и абсолютная интегрируемость.
\item
  Привести пример, где без мажоранты сходимость интегралов ломается.
\end{enumerate}

\subsection{13. Что добавить для
отлично}\label{ux447ux442ux43e-ux434ux43eux431ux430ux432ux438ux442ux44c-ux434ux43bux44f-ux43eux442ux43bux438ux447ux43dux43e}

Для сильного ответа стоит добавить:

\begin{enumerate}
\def\labelenumi{\arabic{enumi}.}
\tightlist
\item
  Схему аппроксимации неотрицательной измеримой функции простыми
  функциями снизу.
\item
  Идею доказательства теоремы Беппо Леви через простую функцию
  \(s\le f\) и множества \(\{f_n\ge\alpha s\}\).
\item
  Идею мажорантной теоремы Лебега: мажоранта не дает массе уходить в
  пики.
\item
  Различие между конечным интегралом, бесконечным интегралом
  неотрицательной функции и интегрируемостью в \(L^1\).
\item
  Связь с математическим ожиданием:
\end{enumerate}

\[
\mathbb E X=\int X\,dP.
\]

\subsection{14. Устная версия ответа на 5
минут}\label{ux443ux441ux442ux43dux430ux44f-ux432ux435ux440ux441ux438ux44f-ux43eux442ux432ux435ux442ux430-ux43dux430-5-ux43cux438ux43dux443ux442}

Интеграл Лебега строится по неотрицательной мере
\((\Omega,\mathcal F,\mu)\). Начинают с индикатора:

\[
\int\mathbf 1_A\,d\mu=\mu(A).
\]

Потом берут простые функции. Простая функция - это измеримая функция,
принимающая конечное число значений. Если

\[
s=\sum_{k=1}^n a_k\mathbf 1_{A_k},
\]

где \(a_k\ge0\), \(A_k\in\mathcal F\) и \(A_k\) попарно не пересекаются,
то

\[
\int s\,d\mu=\sum_{k=1}^n a_k\mu(A_k).
\]

Корректность следует из перехода к общему измельчению двух разбиений.

Для произвольной неотрицательной измеримой функции \(f\) интеграл
определяется как супремум по простым функциям снизу:

\[
\int f\,d\mu
=
\sup_{0\le s\le f,\ s\text{ простая}}
\int s\,d\mu.
\]

Значение может быть бесконечным.

Для знакопеременной функции вводят

\[
f^+=\max(f,0),\qquad f^-=\max(-f,0).
\]

Если

\[
\int |f|\,d\mu<\infty,
\]

то \(f\in L^1\) и

\[
\int f\,d\mu=\int f^+\,d\mu-\int f^-\,d\mu.
\]

Главные теоремы о предельном переходе идут в программе без
доказательства. Теорема Беппо Леви: если

\[
0\le f_n\uparrow f,
\]

то

\[
\int f_n\,d\mu\uparrow\int f\,d\mu.
\]

Мажорантная теорема Лебега: если \(f_n\to f\) почти всюду и есть
\(g\in L^1\) такая, что

\[
|f_n|\le g,
\]

то

\[
\int f_n\,d\mu\to\int f\,d\mu.
\]

Важно: почти всюду сходимости самой по себе недостаточно. Например,
\(f_n=n\mathbf 1_{(0,1/n)}\) сходится к нулю почти всюду на \([0,1]\),
но интегралы равны \(1\).

\subsection{15. Если осталось 2 минуты перед
экзаменом}\label{ux435ux441ux43bux438-ux43eux441ux442ux430ux43bux43eux441ux44c-2-ux43cux438ux43dux443ux442ux44b-ux43fux435ux440ux435ux434-ux44dux43aux437ux430ux43cux435ux43dux43eux43c}

Запомнить:

\[
s=\sum_{k=1}^n a_k\mathbf 1_{A_k},
\qquad
\int s\,d\mu=\sum_{k=1}^n a_k\mu(A_k).
\]

Для \(f\ge0\):

\[
\int f\,d\mu
=
\sup_{0\le s\le f,\ s\text{ простая}}
\int s\,d\mu.
\]

Беппо Леви:

\[
0\le f_n\uparrow f
\Rightarrow
\int f_n\,d\mu\uparrow\int f\,d\mu.
\]

Лебег:

\[
f_n\to f,\quad |f_n|\le g,\quad g\in L^1
\Rightarrow
\int f_n\,d\mu\to\int f\,d\mu.
\]

Главная ловушка: теорема Лебега требует интегрируемую мажоранту, а
просто почти всюду сходимости недостаточно.

\newpage

\end{document}
